在线购买也会同时依赖多个保存域。Alice购买一份组合商品，需要SKU-A和SKU-B各一件。购买权威属于b1，A库存属于b2，B库存属于b3。三个逻辑分区分别由当前副本组承载，各自复制自己的状态。跨分区事务在这些复制组之间，组织一次共同预留。

\CNineteenFigure{ch19-25-participants}{购买协调记录、两份库存子操作和独立平台分别拥有自己的保存位置}{fig:19-participants}

购买分区保存完整$S$、原输入和P/C/T，协调记录把原P关联到稳定事务TX。库存分区保存这个TX针对自己的子操作。支付平台处在另一保存域，由原平台身份关联。图中的横向协议只覆盖内部联合库存，付款沿已建立的C/T协议继续。

\subsection{准备占用进入库存守恒}

先给这次历史的独立初始状态：SKU-A和SKU-B各有4件，都没有准备占用或已提交预留。对每个SKU，$I$表示初始库存，$A$表示当前可用量，$P_{\mathrm{inv}}$表示已经投yes、等待决定的准备占用，$C_{\mathrm{inv}}$表示已提交预留。本节只处理预留，不含补货、发货出库或退回库存，因而保持
\[
I=A+P_{\mathrm{inv}}+C_{\mathrm{inv}}\text{。}
\]
下标inv标明库存计数，区别于购买编号P与付款责任C。

一次成功prepare把一件从$A$移到$P_{\mathrm{inv}}$。commit把它从$P_{\mathrm{inv}}$移到$C_{\mathrm{inv}}$，abort则把它移回$A$。每个动作都在所属库存分区的一次原子转移中完成。已投yes的那件库存如果继续列为可用，其他购买就可能拿走它，先前的准备承诺也会失去资源基础。

\CNineteenFigure{ch19-09}{同一SKU的可用、准备占用与已提交预留保持总量不变}{fig:19-stock}

准备记录要持久保存。参与者投yes后退出，新主或恢复进程从日志及图像取得同一准备占用，等待原TX的决定。若准备只留在旧进程内存，新进程会把那件库存再分配一次，协调者却可能依据先前收到的yes决定commit。

\subsection{稳定TX与不可覆盖的参与者输入}

购买权威受理完整$S$及原输入时分配实际P/C/T。组合商品的C/T先以等待库存资格的状态保存，付款工作者此时无法取得发送资格。协调记录再由原P关联稳定TX，例如\code{TX:P:b1:1}，列出b2、b3及各自的SKU、数量和协调来源。后继协调者读取原记录继续。

参与者按\code{(TX, participant\_id)}保存子操作，内容包括原SKU、数量、协调来源、投票和状态。同一子身份与同样输入到来，返回原投票和当前状态；同一子身份携带不同数量或协调来源，返回冲突并保留原记录。终态为 aborted 时，同参 prepare 返回该终态，库存保持不变。协调者也按参与者身份保存投票，已经保存的yes不能被后来的no覆盖。

参与者投yes前检查可用量并持久形成prepare，库存不足则形成no及拒绝记录。商品组合与参与者清单也绑定到原受理输入。恢复者据此向原参与者重发同一子操作。内部接口的认证和授权限制谁能提交协调命令，业务校验比较清单与商品组合，使另一套资源集合无法覆盖原购买。

\subsection{收票、截止与共同决定}

两阶段提交把准备承诺与共同决定分成两个阶段。协调者收齐全部必要参与者的yes以后才有资格选择commit，收到no时可以选择abort。组合购买还要给缺票的等待设置截止：某个库存分区长期不可达时，已经准备的一件库存需要有释放路径。因此，协调组在首次保存TX时，同时固定参与者清单、原输入和绝对截止时间$d$。后继协调者沿这份记录继续，重试或接管都读取原$d$。

协调状态只有collecting、commit和abort三种。collecting表示决定尚未形成，commit与abort是互斥且不可改写的终态。收齐yes是一项提交条件，决定则由后续的原子命令保存。最后一张yes已经入日志、commit决定尚未入日志时，timeout仍可能先取得abort。决定与超时在同一协调状态机中排序，参与者的本地时钟不参加这项仲裁。

超时由受授权的协调进程或恢复者观察。它读到原截止已到，向原协调组提交携带TX、协调来源、原$d$和观察时间的timeout命令。入口先认证调用者并核验协调权限，状态机再比较这些字段与原TX。各副本应用相同的日志输入，应用命令时不各自读取本机时钟。若业务要求中止绝不早于真实截止，观察者采用具有保守下界的时间源，等下界到达$d$再提交。错误地提前报告时间会使购买过早失败，单一终态的保持则来自状态机的条件转移。

下面将一次命令的原子范围写明。持久提交由协调组的复制协议完成，命令取得提交确认以后，执行者才能把它返回的终态用作通知依据。

\begin{BookCode}
首次建立 TX：
    绑定原 P、参与者清单、各子操作输入和固定截止 d
    state := collecting；votes := 空映射
    重复建立时比较原绑定，冲突则拒绝，匹配则返回原 TX

原子命令 record_vote(TX, participant, vote)：
    核验参与者属于原清单，vote 为 yes 或 no
    若 state != collecting：返回原终态
    若已有不同投票：返回 vote_conflict
    保存该参与者的原投票；返回当前 TX

原子命令 decide(TX)：
    若 state != collecting：返回原终态
    若原清单中有 no：state := abort
    否则若原清单已收齐 yes：state := commit
    否则：返回 collecting
    保存并返回该终态

原子命令 timeout(TX, coordinator, expected_d, observed_time)：
    核验协调来源、原截止及观察输入
    若 state != collecting：返回原终态
    若 observed_time < d：返回 deadline_not_reached
    state := abort；保存并返回该终态

复制提交确认后：
    若返回 collecting：继续收票或等待截止观察
    若返回 commit 或 abort：按原清单反复传播同一决定
    回复未知时先读原协调记录，再决定下一动作
\end{BookCode}

三种次序给出竞态的不同结果。最后yes、decide、timeout依次应用时，decide先保存commit，timeout读回commit。最后yes、timeout、decide依次应用时，timeout先保存abort，即使全部投票已经齐备，decide也读回abort。timeout先于最后yes时，迟票直接取得abort，原投票集合保持终态形成时的内容。两项终态命令都先检查collecting，所以每一种允许迁移都保持终态唯一。

\CNineteenFigure{ch19-38-timeout-arbitration}{同一协调日志串行仲裁决定与超时，迟票及重试读取已固定终态}{fig:19-timeout}

这里的$d$表示协调者取得中止资格的时点。它到达以后，若timeout尚未进入协调序列，具备全部yes的decide依然可能先形成commit。需要按严格业务截止拒绝提交的系统，会把对应的可信时间观察及拒绝条件也纳入decide命令。采用本节规则时，客户端查询TX的已提交终态，界面倒计时本身没有决定效力。

\subsection{先持久决定，再传播}

\CNineteenFigure{ch19-10}{参与者先保存准备承诺，协调组持久保存共同决定，再传播到各参与者}{fig:19-2pc}

commit或abort先在协调权威组取得持久复制，然后才通知参与者。同一个TX的决定形成后保持不变。通知带上原TX、参与者身份、协调来源和原SKU及数量。参与者收到commit，把原准备占用变为已提交预留；收到abort，释放原占用。重复通知比较完整关联并返回原终态，commit之后的abort或abort之后的commit进入决定冲突。

决定来源通过明确接口核验。部署可以让参与者查询协调权威的已提交TX记录，也可以让受控恢复服务先读权威，再发送内部通知。后一结构中，参与者只接受被授权的服务调用，服务承担读取和核对原决定的责任。任意客户端提供一个\code{finish(commit)}参数，不能取得资源提交资格。

持久决定和各参与者的应用进度可能不同步。b2已经commit、b3尚处于prepared时，共同决定可以已经固定为commit。查询如实显示推进中，恢复者把同一决定传到b3。尚未收到通知的参与者继续保存准备，等待的时间长短不改变原决定。

\subsection{丢票、丢通知与恢复责任}

b2把可用量从4改成3，并保存一件准备占用，yes回复在网络中丢失。协调者重发原子操作身份，b2读取原记录，返回yes。若协调者改造一个新TX，b2会将其作为另一项准备处理，第一项责任则继续悬挂。

协调者也可能已持久保存commit，却在通知前整个组不可达。对一个只持有prepared的参与者而言，本地观察同时兼容两条历史：协调者选了commit，或者另一参与者投no、协调者选了abort。它要根据可核验的共同决定推进，单靠等待时间无法区分这两条历史。

\CNineteenFigure{ch19-11}{协调决定暂不可读时保留准备占用，恢复后重发原决定}{fig:19-2pc-recovery}

协调组恢复多数以后，后继工作者读取TX的原状态。已经commit或abort，就向原参与者重发该决定。处于collecting且原截止已到，就重新提交timeout；截止尚未到时，按原参与者清单继续收票。b2已经commit而回复丢失时，再次通知返回原已提交预留。b2的主进程在prepared时退出，则由新主从多数复制中恢复同一准备，再接受原决定。固定截止、投票与终态随协调日志及图像保存，恢复者无需从旧进程的定时器推断结果。

两阶段提交的阻塞也在这里出现：协调组失去多数时，timeout也无法持久形成abort，准备资源继续占用。协调组能够提交以后，截止观察、恢复者调度和通知送达共同推动释放；它们决定进展，单一终态则由前面的条件转移保持。运维可以限制新组合购买，优先恢复现有TX。参与者若仅因本地等待超时就删除prepare，会在原决定为commit的历史中释放已经承诺的资源。

abort分支保持同样的身份关联。另取一条独立历史：A可用4件，B可用0件，TX要求各一件。b2成功prepare，b3保存no，协调者持久保存abort。b2将原准备的一件移回可用，只释放一次。abort回复丢失后重复通知，A保持4件，反向commit进入冲突。

同一初态下，给 b2 的首次 prepare 也可能滞留，b3 的 no 或协调者的timeout先使TX进入abort。b2收到可信abort时尚无子操作，它在一次原子转移中保存aborted终态、原协调来源和输入，可用量保持4。迟到的同参prepare查到该终态，返回aborted；异参prepare返回冲突。协调者记录b2完成以后，迟到请求也就无法重新产生prepared占用。

本版把这些终态随参与者图像和日志保存，迁移时一并带走。要回收终态，系统必须先形成持久关闭边界，使旧 TX 的 prepare 和 finish 在所有合法入口都被拒绝，再释放相应记录。单靠协调者已经结束或等待足够久，无法阻止网络中的旧 prepare 后到。

\CNineteenFigure{ch19-26-abort}{abort覆盖已有准备和首次prepare尚未到达两种状态，终态阻止重新占用}{fig:19-abort}

\subsection{库存决定授予原C/T付款资格}

内部联合预留的持久commit是所有实际付款发送路径的必经条件。C/T已经保存，并不意味着可以调用平台。本书的参考流程先确认原TX的commit，再推进参与者并记录完成，随后把原C/T的\code{reservation\_ready}置为true。支付扫描器、人工恢复入口和迟到工作者的发送路径都核验这个条件。

\CNineteenFigure{ch19-27-payment-gate}{先保存原C/T，库存责任闭合后释放发送资格，付款继续使用原平台身份}{fig:19-payment-gate}

本节付款金额为3000分，平台初始余额为10000分。取得资格以后，支付工作者用原scope和\code{pay-C}发送。工作者重启、库存恢复或客户端重试都沿同一$S$、P/C/T与TX继续。只在页面隐藏支付按钮无法约束后台扫描和迟到请求，资格核验位于实际发送链路。

库存commit后、付款结果尚未确定时，库存保持已预留。付款超时后查询原平台操作，已有成功便恢复原结果。确定拒绝、订单取消或履约变化，由后续业务流程发起释放或其他反向动作，并关联原预留。它们是原事务提交以后的新转移，原TX的commit作为历史事实保留。

\subsection{独立付款与分阶段业务}

普通商业支付HTTP接口通常没有承诺进入本系统prepare并持有资源等待共同决定。把调用写在协调函数里，平台账本也不会自动成为库存事务参与者。付款通过稳定平台身份、查询与对账，衔接两个保存域各自形成的事实。

若业务允许阶段性可见，可以采用Saga流程：先预留A，再预留B，再付款，失败后按已完成阶段执行补偿。每个阶段保存原身份、输入、结果及下一责任。A已经预留、B尚未完成时，这个中间状态进入查询与库存可用性规则，恢复者据其继续~\cite{sagas}。

\CNineteenFigure{ch19-12}{Saga保存已完成阶段，并以独立身份继续释放或退款}{fig:19-saga}

释放预留使用稳定释放身份，关联原预留；退款使用新的RF，关联原C与平台收款结果。补偿超时以后按原补偿身份查询或重试。原付款成功继续保留，退款可能处于处理中，成功与未知退款共同占用可退额度。完整图像保存这些关联，后继执行者才能继续同一补偿。

内部联合库存要求共同预留时，两阶段提交固定一个共同决定。允许部分阶段先行可见时，业务状态及补偿责任给出另一种组织。每种安排都把等待占用、可见阶段和恢复负责人写入保存结构，后续付款和取消据这些事实推进。
\FloatBarrier
